\section{Умножение}

\begin{definition}{2.3.1}{умножение натуральных чисел}
Пусть $m$ — натуральное число. Умножение нуля на $m$ определяем как $0 \times m := 0$.
Предположим по индукции, что умножение $n$ на $m$ уже определено. Тогда умножение
$n{+}{+}$ на $m$ определяем как $(n{+}{+}) \times m := (n \times m) + m$.
\leansource{Nat.mul}{(n m : Nat) : Nat}
\end{definition}

\begin{remark}
Умножение — это итерированное сложение, как сложение — итерированное взятие следующего
числа. Дальше $n \times m$ будем сокращённо писать как $nm$, с обычным приоритетом
умножения над сложением: $ab+c$ значит $(a \times b)+c$. Скобки нужны, чтобы записать
$a \times (b+c)$.
\end{remark}

\begin{lemma}{2.3.2}{умножение коммутативно}
Для любых натуральных $n,m$ выполняется $nm=mn$.
\leansource{Nat.mul_comm}{(n m : Nat) : n * m = m * n}
\end{lemma}

\begin{proof}
Индукция по $n$ при фиксированном $m$, с двумя вспомогательными фактами: $n \times 0=0$
(индукцией по $n$, аналогично \taoref{lem:2.2.2}{лемме 2.2.2}) и
$n \times (m{+}{+})=nm+n$ (индукцией по $n$, аналогично \taoref{lem:2.2.3}{лемме 2.2.3}).\\
\prooflabel{База} $0 \times m=0$ по определению, и $m \times 0=0$ по первому вспомогательному
факту.\\
\prooflabel{Переход} Пусть $nm=mn$. Тогда $(n{+}{+})m=nm+m$ по определению, а
$m(n{+}{+})=mn+m$ по второму вспомогательному факту — обе части равны, поскольку $nm=mn$.
\end{proof}

\begin{lemma}{2.3.3}{у положительных чисел нет делителей нуля}
Пусть $n,m$ — натуральные числа. Тогда $nm=0$ равносильно тому, что $n=0$ или $m=0$. В
частности, если оба положительны, то и $nm$ положительно.
\leansource{Nat.mul_eq_zero}{(n m : Nat) : n * m = 0 ↔ n = 0 ∨ m = 0}
\end{lemma}

\begin{remark}
В одну сторону — прямая проверка по \taoref{def:2.3.1}{определению}: если $n=0$ или $m=0$,
то $nm=0$. В другую — от противного, индукцией по \taoref{prop:2.2.8}{утверждению 2.2.8}
(сумма положительного и натурального положительна).
\end{remark}

\begin{proposition}{2.3.4}{дистрибутивный закон}
Для любых натуральных $a,b,c$ выполняется $a(b+c)=ab+ac$ и $(b+c)a=ba+ca$.
\leansource{Nat.mul_add}{(a b c : Nat) : a * (b + c) = a * b + a * c}
\leansource{Nat.add_mul}{(a b c : Nat) : (a + b) * c = a * c + b * c}
\end{proposition}

\begin{proof}
По \taoref{lem:2.3.2}{коммутативности} достаточно доказать первое равенство. Индукция по
$c$ при фиксированных $a,b$.\\
\prooflabel{База} $a(b+0)=ab$ по \taoref{lem:2.2.2}{лемме 2.2.2}, и $ab+a \times 0=ab+0=ab$
— обе части равны $ab$.\\
\prooflabel{Переход} Пусть $a(b+c)=ab+ac$. Тогда $a(b+c{+}{+})=a((b+c){+}{+})=a(b+c)+a$ по
определению умножения, что по предположению индукции равно $(ab+ac)+a=ab+(ac+a)=ab+ac{+}{+}$.
\end{proof}

\begin{proposition}{2.3.5}{умножение ассоциативно}
Для любых натуральных $a,b,c$ выполняется $(ab)c=a(bc)$.
\leansource{Nat.mul_assoc}{(a b c : Nat) : (a * b) * c = a * (b * c)}
\end{proposition}

\begin{proof}
Индукция по $a$ при фиксированных $b,c$.\\
\prooflabel{База} $(0 \times b)c = 0 \times c = 0$ и $0 \times (bc) = 0$ — обе части равны
нулю.\\
\prooflabel{Переход} Пусть $(ab)c=a(bc)$. Тогда $((a{+}{+})b)c=(ab+b)c=(ab)c+bc$ по
\taoref{prop:2.3.4}{дистрибутивности}, что по предположению индукции равно
$a(bc)+bc=(a{+}{+})(bc)$ — снова по дистрибутивности.
\end{proof}

\begin{proposition}{2.3.6}{умножение сохраняет порядок}
Пусть $a,b$ — натуральные числа, $a<b$, и $c$ положительно. Тогда $ac<bc$.
\leansource{Nat.mul_lt_mul_of_pos_right}{(h : a < b) (hc : c.IsPos) : a * c < b * c}
\end{proposition}

\begin{proof}
Из $a<b$ по \taoref{prop:2.2.12}{утверждению 2.2.12(f)} следует $b=a+d$ для некоторого
положительного $d$.\by{\taoref{prop:2.2.12}{утверждение 2.2.12(f)}}\\
Умножая обе части на $c$ и раскрывая дистрибутивность, получаем $bc=ac+dc$.\by{\taoref{prop:2.3.4}{утверждение 2.3.4}}\\
$dc$ положительно, поскольку $d$ и $c$ положительны.\by{\taoref{lem:2.3.3}{лемма 2.3.3}}\\
Значит $ac<bc$.\by{\taoref{prop:2.2.12}{утверждение 2.2.12(f)}}
\end{proof}

\begin{corollary}{2.3.7}{закон сокращения}
Пусть $a,b,c$ — натуральные числа, $ac=bc$ и $c$ положительно. Тогда $a=b$.
\leansource{Nat.mul_cancel_right}{(h : a * c = b * c) (hc : c.IsPos) : a = b}
\end{corollary}

\begin{proof}
По \taoref{prop:2.2.13}{трихотомии} возможны три случая: $a<b$, $a=b$, $a>b$.\\
\caselabel{$a<b$} Тогда $ac<bc$ по \taoref{prop:2.3.6}{утверждению 2.3.6}, что противоречит
$ac=bc$.\\
\caselabel{$a>b$} Тогда $ac>bc$ по тому же утверждению — снова противоречие.\\
Остаётся единственная возможность $a=b$.
\end{proof}

\begin{remark}
Как и \taoref{prop:2.2.6}{закон сокращения для сложения}, это «виртуальное деление»: оно
позволяет сокращать общий положительный множитель, ещё не имея деления как отдельной
операции.
\end{remark}

\begin{proposition}{2.3.9}{алгоритм деления с остатком}
Пусть $n$ — натуральное число, $q$ — положительное. Тогда существуют натуральные $m,r$
такие, что $0 \leq r < q$ и $n=mq+r$.
\leansource{Nat.exists_div_mod}{(n : Nat) (hq : q.IsPos) : ∃ m r : Nat, 0 ≤ r ∧ r < q ∧ n = m * q + r}
\end{proposition}

\begin{remark}
Это деление $n$ на $q$ с частным $m$ и остатком $r$. Доказательство (индукцией по $n$ при
фиксированном $q$) здесь опущено — оно есть в Lean-файле, а идею стоит продумать
самостоятельно, как и в книге.
\end{remark}

\begin{definition}{2.3.11}{возведение в степень}
Пусть $m$ — натуральное число. Возведение $m$ в степень $0$ определяем как $m^0 := 1$
(в частности, $0^0 := 1$). Предположим по индукции, что $m^n$ уже определено. Тогда
$m^{n{+}{+}} := m^n \times m$.
\leansource{Nat.pow}{(m n : Nat) : Nat}
\end{definition}

\begin{remark}
Как умножение — итерированное сложение, возведение в степень — итерированное умножение:
$m^1=m$, $m^2=m \times m$, $m^3=m \times m \times m$ и так далее.
\end{remark}
